\begin{description}
\item[a.] The center of gravity of the stars is relatively to the star A given
  by
  \[ r_A = \frac{m_B}{m_A + m_B}\,d \]
  and since the orbit of A is a circle, then
  \[ F_{AX} = \frac{G m_A m_B}{d^2} = \frac{m_A v_A^2}{r_A} \]
  So, we get
  \[ v_A = m_B\sqrt{\frac{G}{(m_A + m_B)d}} \]
  The angular velocity of A is given by
  \[ \omega = \frac{v_A}{r_A} = \sqrt{\frac{G(m_A + m_B)}{d^3}} \]
  \hfill(30)
\item[b.] In Cartesian coordinate system, the velocity of A is
  \[ \vec{v}_A(t) = v_A\left(-\sin(\omega t + \varepsilon)\,\hat{\imath}
     + \cos(\omega t + \varepsilon)\,\hat{\jmath}\right) \]
  Unit vector of the observer is
  \[ \hat{r}_P = \cos\theta\,\hat{k} + \sin\theta\,\hat{\jmath} \]
  so the component of $\vec{v}_A$ in the line of the observer sight is given
  by
  \[ \vec{v}_A\cdot\hat{r}_P = v_A\sin\theta\,\cos(\omega t + \varepsilon) \]
  Since the component of $\vec{v}_A$ in the line of the observer sight is
  $K\cos(\omega t + \varepsilon)$, then
  \[ K = v_A\sin\theta \]
  Finally, we have
  \[ \frac{K^3}{\omega G} = \frac{m_B^3}{(m_A + m_B)^2}\sin^3\theta
     \qquad (**) \]
  \hfill(30)
\item[c.] From the result in b., namely eq.\ (**), we get
  \[ \sin^3\theta = \frac{K^3}{\omega G}\,
     \frac{(m_A + m_B)^2}{m_B^3}
     < \frac{1}{250}\,\frac{32^2}{8} = \frac{64}{125} \]
  Since $\theta \in [0, \pi]$, then $\sin\theta < 0.8$. Thus, the probability
  of B is a black hole is the same as the probability of $\sin\theta < 0.8$
  for $\theta \in [0, \pi]$. So $\theta$ is less than $53^\circ$ or greater
  than $127^\circ$.
  \hfill(40)
\end{description}
